\section{Architektur-Grundlagen \& Frameworks}

\subsection{Warum wichtig?}
Pläne beeinflussen spätere Entscheidungen. Viele Änderungsanfragen
sind vermeidbar, sonst entstehen \emph{technische Schulden}. Framework-Änderungen
allein lösen keine Grundprobleme, und eingeschränkte Komponentenwahl
wirkt sich auf Folgeprojekte aus (z.\,B. Lizenzen).

\subsection{Technische Schulden}
\textbf{Definition:} Konsequenzen bei schlechter Umsetzung.\\[4pt]
\textbf{Arten:}
\begin{itemize}
    \item Bewusste Schulden (strategische Entscheidungen)
    \item Unbeabsichtigte Schulden durch mangelnde Dokumentation,
          fehlende Automatisierung, fehlende Coding Standards
          oder Missachtung von Code-Hinweisen
\end{itemize}

\subsection{Ziele einer Lösungs-Architektur}
\begin{itemize}
    \item Skalierbarkeit, Resilienz, Sicherheit, Compliance
    \item Globaler Kundendienst mit gleicher Performance
    \item Wartbarkeit über Jahrzehnte
\end{itemize}

\subsection{Architektur-Perspektiven}
\begin{itemize}
    \item \textbf{Geschäftssicht:} Ziele, Vision, Stakeholder,
          Geschäftsanforderungen, Erfolgsmetriken
    \item \textbf{Prozesssicht:} Abläufe von Kunde zu Server,
          z.\,B. als Sequenzdiagramm oder BPMN
    \item \textbf{Konzeptionelle Sicht:} Architektur ohne
          spezifische Cloud-Provider-Komponenten
    \item \textbf{Lösungssicht:} Cloud-Service-Architektur,
          User Flow, ER-Diagramm
    \item \textbf{Sicherheitssicht:} Zugangsverwaltung,
          Compliance-Anforderungen, Firewall
\end{itemize}

\newpage

\subsection{Architektur-Frameworks: Governance}
Dienen der unternehmensweiten Steuerung, Compliance-Nachweisen
und Risikokontrolle über vordefinierte Prinzipien, Gremien und Reviews.
Beispiele: \textit{TOGAF}, \textit{ISO/IEC/IEEE 42010:2022}, \textit{Zachman-Framework}.\\[4pt]
\textbf{TOGAF} (The Open Group Architecture Framework): Framework zur
Entwicklung und Verwaltung von Unternehmensarchitekturen (People,
Processes, Technology) mit vier Architektur-Domänen:
% LEKTORAT (fachlich): Vorher als "vier Domänen" angekündigt, aber nur drei
% gelistet. TOGAF unterscheidet vier Domänen; die Informationssystem-
% architektur zerfällt in Anwendungs- und Datenarchitektur.
\begin{itemize}
    \item Geschäftsarchitektur (Business)
    \item Datenarchitektur (Data)
    \item Anwendungsarchitektur (Application)
    \item Technologiearchitektur (Technology)
\end{itemize}

\subsection{Architektur-Frameworks: Praxisorientiert}
Stellen schnelle, verständliche Diagramme für Entwickler, DevOps und
Produktteams bereit. Zielen eher auf einzelne Lösungen und Applikationen
ab. Beispiele: \textit{C4-Modell}, \textit{4+1-Modell}, \textit{PlantUML}.\\[4pt]
\textbf{C4-Modell:} Unterstützt Software-Entwicklungsteams, ihre
Architektur bei Planung und Durchführung zu erklären.
4 Ebenen: Software System, Container, Component, Code.\\[4pt]
\textbf{Software Architecture Diagramming Maturity Model:}
\begin{itemize}
    \item \textbf{Level 1 -- Initial:} Keine Architekturdiagramme
    \item \textbf{Level 2 -- Ad hoc:} Diagramme mit ad-hoc-Abstraktionen
          in einem allgemeinen Diagramm-Tool
    \item \textbf{Level 3 -- Defined:} Diagramme mit definierten
          Abstraktionen und Notation, allgemeines Tool
    \item \textbf{Level 4 -- Modelled:} Definierte Abstraktionen in
          einem Modellierungs-Tool, manuell erstellt
    \item \textbf{Level 5 -- Optimising:} Modellelemente werden
          teamübergreifend geteilt, als querybare Datensätze genutzt
          und automatisch aus Sourcecode bzw. Deployment-Umgebung
          generiert
\end{itemize}

\subsection{Frameworks im Detail (TOGAF \& C4)}
\subsubsection{TOGAF -- Architecture Development Method (ADM)}
\begin{figure}[H]
    \centering
    \includegraphics[width=\textwidth]{TOGAF.png}
    \caption{TOGAF ADM-Zyklus}
\end{figure}

Der ADM-Zyklus beschreibt, wie eine Unternehmensarchitektur schrittweise
entwickelt wird. Im Zentrum stehen immer die \textbf{Requirements} (Anforderungen),
auf die jede Phase kontinuierlich einzahlt.

\begin{itemize}
    \item \textbf{Preliminary:} Framework und Prinzipien festlegen
          (Vorbereitung, noch kein offizieller Buchstabe)
    \item \textbf{A -- Architecture Vision:} Ziele, Scope und
          Stakeholder-Erwartungen definieren
    \item \textbf{B -- Business Architecture:} Geschäftsprozesse
          und -strukturen modellieren
    \item \textbf{C -- Information System Architecture:}
          Anwendungs- und Datenarchitektur beschreiben
    \item \textbf{D -- Technology Architecture:} Infrastruktur,
          Plattformen und Technologien festlegen
    \item \textbf{E -- Opportunities \& Solutions:} Lücken zwischen
          Ist- und Soll-Zustand identifizieren, Lösungsoptionen bewerten
    \item \textbf{F -- Migration Planning:} Roadmap und Reihenfolge
          der Umsetzung planen
    \item \textbf{G -- Implementation Governance:} Umsetzung
          begleiten und sicherstellen, dass Architekturvorgaben
          eingehalten werden
    \item \textbf{H -- Architecture Change Management:} Änderungen
          am laufenden System kontrolliert einführen und den
          Zyklus bei Bedarf neu starten
\end{itemize}

Der Zyklus ist \textbf{iterativ}, d.\,h. nach Phase H kann der Prozess
wieder bei A (oder sogar Preliminary) beginnen, falls sich Anforderungen
grundlegend ändern.

\subsubsection{C4}

Das C4-Modell ist ein praxisorientiertes Framework zur Visualisierung
von Software-Architektur. Der Name steht für die vier Abstraktionsebenen:
\textbf{Context, Container, Component und Code}. Die Idee dahinter ist,
dass verschiedene Zielgruppen unterschiedliche Detailtiefe brauchen,
ein Manager braucht nur den groben Überblick, ein Entwickler hingegen
die genaue Komponentenstruktur.
\begin{figure}[H]
    \centering
    \includegraphics[width=\textwidth]{c4.png}
    \caption{C4-Modell -- vier Abstraktionsebenen}
\end{figure}

Die vier Ebenen bauen aufeinander auf:
\begin{itemize}
    \item \textbf{System Context:} Zeigt das System als Ganzes
          und seine Beziehungen zu Nutzern und externen Systemen
    \item \textbf{Container:} Aufschlüsselung in deploybare Einheiten
          wie Web-Apps, Datenbanken oder Microservices
    \item \textbf{Component:} Interne Struktur eines Containers,
          z.\,B. einzelne Module oder Services
    \item \textbf{Code:} Klassendiagramme oder ähnliches auf
          Code-Ebene (wird selten gezeichnet, da IDEs das übernehmen)
\end{itemize}

Ein großer Vorteil des C4-Modells ist die \textbf{Werkzeugfreiheit}.
Es funktioniert mit PlantUML, draw.io oder sogar auf einem Whiteboard,
solange die vier Ebenen klar getrennt bleiben.
